\color{unidarkgreen}{\subsection{Реле тока блокировки РПН}}
\color{black}

\begin{enumerate}[label=\arabic{section}.\arabic{subsection}.\arabic*,  align=left, labelsep=4pt, leftmargin=0pt, itemindent=57pt]

\item
В целях блокирования автоматики РПН при превышении током через РПН допустимого значения предусматривается реле блокировки. По умолчанию данное реле подключено к трансформаторам тока стороны СН автотрансформатора. Порог по току срабатывания реле задается параметром \textbf{«Iбл РПН»}.

\item
Алгоритм работы РТ РПН приведен на рисунке \ref{pic:rtrpn}. Уставки РТ РПН приведены в таблице \ref{app_set:rtrpn}.

\vspace{2mm}
\begin{figure}[H]
\centering
\includegraphics[width=0.9\textwidth,height=0.9\textheight,keepaspectratio]{\fbpath/041.1181 РПН/img/RPN.pdf}
\captionof{figure}{Функциональная схема логики РТ блокировки РПН}\label{pic:rtrpn}
\end{figure}

\end{enumerate}